% TODO make hw stylesheet
\documentclass[10pt]{article}
\usepackage[letterpaper,top=1in]{geometry}
\usepackage[dvipsnames,svgnames]{xcolor}
\usepackage[shortlabels]{enumitem}
\usepackage[framemethod=TikZ]{mdframed}
\usepackage{amsmath,amssymb,amsthm}
\usepackage[colorlinks]{hyperref}
\usepackage{multicol}
\usepackage{tabularx}

\hypersetup{
	colorlinks=true,
	linkcolor=blue,
	filecolor=magenta,
	urlcolor=magenta,
	pdftitle={MAT 528 HW}
}

\usepackage{mathtools}
\usepackage{thmtools}
\usepackage[textsize=footnotesize]{todonotes}
\usepackage{titling}

\reversemarginpar
\mdfdefinestyle{mdbluebox}{roundcorner=10pt,innerbottommargin=9pt,
	linecolor=blue,backgroundcolor=TealBlue!5,}
\declaretheoremstyle[headfont=\sffamily\bfseries\color{MidnightBlue},
mdframed={style=mdbluebox},]{thmbluebox}

\mdfdefinestyle{mdbloobox}{frametitlefont=\bfseries,innerbottommargin=8pt,
	nobreak=true,backgroundcolor=TealBlue!5,linecolor=blue,}
\declaretheoremstyle[headfont=\bfseries\color{MidnightBlue},
mdframed={style=mdbloobox},headpunct={\\[3pt]},postheadspace=0pt,]{thmbloobox}

\mdfdefinestyle{mdredbox}{frametitlefont=\bfseries,innerbottommargin=8pt,
	nobreak=true,backgroundcolor=Salmon!5,linecolor=RawSienna,}
\declaretheoremstyle[headfont=\bfseries\color{RawSienna},
mdframed={style=mdredbox},headpunct={\\[3pt]},postheadspace=0pt,]{thmredbox}

\mdfdefinestyle{mdgreenbox}{linecolor=ForestGreen,backgroundcolor=ForestGreen!5,
	linewidth=2pt,rightline=false,leftline=true,topline=false,bottomline=false,}
\declaretheoremstyle[headfont=\bfseries\sffamily\color{ForestGreen!70!black},
mdframed={style=mdgreenbox},headpunct={ --- },]{thmgreenbox}

\mdfdefinestyle{mdorangebox}{linecolor=RedOrange,backgroundcolor=RedOrange!5,
	linewidth=2pt,rightline=false,leftline=true,topline=false,bottomline=false,}
\declaretheoremstyle[headfont=\bfseries\sffamily\color{RedOrange!70!black},
mdframed={style=mdorangebox},headpunct={ --- },]{thmorangebox}

\mdfdefinestyle{mdblackbox}{linecolor=black,backgroundcolor=RedViolet!5!gray!5,
	linewidth=3pt,rightline=false,leftline=true,topline=false,bottomline=false,}
\declaretheoremstyle[mdframed={style=mdblackbox}]{thmblackbox}

\declaretheoremstyle[spaceabove=6pt,spacebelow=6pt]{basehead}

\declaretheorem[style=basehead,name=Problem]{problem}
\declaretheorem[style=thmredbox,name=Problem,numbered=no]{problem*}
\declaretheorem[style=thmbloobox,name=Setup]{setup} % for config geo
\declaretheorem[style=thmredbox,name=Example,sibling=problem]{example}
\declaretheorem[style=thmredbox,name=$\bigstar$ Problem,sibling=problem]{niceprob}
\declaretheorem[style=thmbluebox,name=Theorem,sibling=problem]{theorem}
\declaretheorem[style=thmbluebox,name=Corollary,sibling=theorem]{corollary}
\declaretheorem[style=thmbluebox,name=Proposition,sibling=theorem]{prop}
\declaretheorem[style=thmbluebox,name=Lemma,numberwithin=problem]{lemma}
\declaretheorem[style=thmbluebox,name=Theorem,numbered=no]{theorem*}
\declaretheorem[style=thmbluebox,name=Lemma,numbered=no]{lemma*}
\declaretheorem[style=thmgreenbox,name=Claim,numberwithin=problem]{claim}
\declaretheorem[style=thmgreenbox,name=Claim,numbered=no]{claim*}
\declaretheorem[style=thmblackbox,name=Remark,numberwithin=problem]{remark}
\declaretheorem[style=thmblackbox,name=Remark,numbered=no]{remark*}
\declaretheorem[style=thmgreenbox,name=Definition,sibling=theorem]{definition}
\declaretheorem[style=thmgreenbox,name=Definition,numbered=no]{definition*}

\providecommand{\ol}{\overline}
\providecommand{\ceil}[1]{\left\lceil #1 \right\rceil}
\providecommand{\floor}[1]{\left\lfloor #1 \right\rfloor}
\newcommand{\dd}{\mathrm{d}}
\providecommand{\ul}{\underline}
\providecommand{\eps}{\varepsilon}
\providecommand{\half}{\frac{1}{2}}
\providecommand{\CC}{\mathbb C}
\providecommand{\FF}{\mathbb F}
\providecommand{\NN}{\mathbb N}
\providecommand{\QQ}{\mathbb Q}
\providecommand{\RR}{\mathbb R}
\providecommand{\ZZ}{\mathbb Z}
\providecommand{\dg}{^\circ}
\providecommand{\ii}{\item}
\providecommand{\setdiff}{\smallsetminus}
\providecommand{\alert}[1]{{\sffamily\textbf{\textcolor{blue}{#1}}}}
\providecommand{\inv}{^{-1}}
\providecommand{\mb}[1]{\mathbf{#1}}

\newenvironment{soln}{\begin{proof}[Solution]}{\end{proof}}

\providecommand{\pow}{\operatorname{Pow}}
\providecommand{\vocab}[1]{{\textbf{\textcolor{ForestGreen}{#1}}}}
\providecommand{\lcm}{\operatorname{lcm}}
\providecommand{\abs}[1]{\left\lvert #1\right\rvert}
\providecommand{\norm}[1]{\left\|#1\right\|}

\usepackage{mathrsfs}

% place title at top right
\setlength{\droptitle}{-8ex}
\pretitle{\begin{flushright}\Large\bfseries}
\posttitle{\par\end{flushright}}
\preauthor{\begin{flushright}}
\postauthor{\end{flushright}}
\predate{\begin{flushright}}
\postdate{\end{flushright}}

\title{MAT 528 HW12}
\author{\large Aditya Pahuja}
\date{\large Due 12/08}
\begin{document}
\maketitle

\begin{problem}
    [Munkres 58.2b]
    Determine whether the fundamental group of
    the torus $T$ with a point removed
    is trivial, infinite cyclic, or isomorphic to
    the fundamental group of the figure eight.
\end{problem}

\begin{soln}
    Let $T'$ be $T$ with a point removed.
    Recall that $T$ can be obtained by identifying opposite edges
    of a square, so $T'$ can be obtained by identifying opposite edges of
    a square with a point $x$ removed from its interior.
    This square-minus-a-point deformation retracts onto its boundary
    (specifically, for each $p$ in the interior of the square,
    let $q$ be the intersection of the ray from $x$ through $p$
    with the boundary of the square, and move $p$ linearly along
    the segment from $p$ to $q$),
    so the image of the boundary under the quotient map
    identifying the edges of the square
    is a deformation retract of $T'$. On the other hand
    the boundary with its edges identified is simply the figure eight,
    so $T'$ has the same fundamental group
    as the figure eight.
\end{soln}

\begin{problem}
    [Munkres 58.5]
    Recall that a space $X$ is said to be \vocab{contractible} if
    the identity map from $X$ to itself is nullhomotopic.
    Show that $X$ is contractible if and only if $X$ has
    the homotopy type of a one-point space.
\end{problem}

\begin{soln}
    Suppose first that $X$ is contractible.
    Then there is a homotopy $H\colon X\times I\to X$
    such that $H(x,0) = x$ and $H(x,1) =x_0$ for some fixed $x_0\in X$
    and all $x\in X$. Let $i$ be the inclusion map $i\colon\{x_0\}\to X$
    and let $f\colon X\to\{x_0\}$ be the map $f(x) = x_0$.
    Then $f\circ i\colon\{x_0\}\to\{x_0\}$ is (homotopic to) the identity on $\{x_0\}$
    and $i\circ f\colon X\to X$ sends every point in $x$ to $x_0$,
    which via the homotopy $H$ is homotopic to the identity on $X$;
    this proves that $f$ and $i$ are homotopy equivalences between $X$
    and a one-point space.

    Conversely, suppose $X$ has the homotopy type of a one-point space $Y$
    and let $f\colon X\to Y$, $g\colon Y\to X$ be homotopy equivalences.
    Then $g\circ f\colon X\to X$ is the constant map sending
    every point in $X$ to the unique point in the image of $g$
    and is homotopic to the identity on $X$ by the assumption
    that $f$ and $g$ are homotopy equivalences.
    Thus the identity on $X$ is nullhomotopic, i.e. $X$ is contractible.
\end{soln}

\begin{problem}
    [Munkres 58.9abcd]
    Let $b_0$ be the point $(1,0)\in S^1$;
    choose a generator $\gamma$ for the infinite cyclic group $\pi_1(S^1,b_0)$.
    If $x_0\in S^1$, choose a path $\alpha\colon[0,1]\to S^1$ from $b_0$ to $x_0$
    and define $\gamma_{x_0} = \widehat\alpha(\gamma)$.
    Then $\gamma_{x_0}$ generates $\pi_1(S^1,x_0)$.
    The element $\widehat\alpha(\gamma)$ is independent
    of the choice of the path $\alpha$,
    since the fundamental group of $S^1$ is abelian.

    Now given $h\colon S^1\to S^1$, choose $x_0\in S^1$ and let $h(x_0) = x_1$.
    Consider the homomorphism $h_*\colon\pi_1(S^1,x_0)\to\pi_1(S^1,x_1)$.
    Since both groups are infinite cyclic,
    $h_*(\gamma_{x_0}) = d\cdot\gamma_{x_1}$ for some integer $d$
    (writing the group operation additively);
    this $d$ is called the \vocab{degree} of $h$,
    denoted $\deg h$.
    \begin{enumerate}[(a)]
        \ii Show that $d$ is independent of the choice of $x_0$.
	\ii Show that if $h,k\colon S^1\to S^1$ are homotopic,
	they have the same degree.
	\ii Show that $\deg(h\circ k) = (\deg h)(\deg k)$.
	\ii Compute the degrees of the constant map, the identity map,
	the reflection map $\rho(x_1,x_2) = (x_1,-x_2)$,
	and the map $h(z) = z^n$, where $z$ is a complex number.
    \end{enumerate}
\end{problem}

\begin{soln}
    (a) Let $y_0\in S^1$ and $y_1 = h(y_0)$.
    Let $\alpha_j$ be a path from $x_j$ to $y_j$ for $j=1,2$.
    Let $h_{*x}\colon\pi_1(S^1,x_0)\to\pi_1(S^1,x_1)$,
    $h_{*y}\colon\pi_1(S^1,y_0)\to\pi_1(S^1,y_1)$ be
    the induced homomorphisms by $h$.
    I claim that $\widehat{\alpha_1}\circ h_{*x} = h_{*y}\circ\widehat{\alpha_0}$.
    To see this, let $g_{x_0}$ be a path in the path homotopy class $\gamma_{x_0}$.
    We have
    \begin{equation*}
	h_{*y}(\widehat{\alpha_0}([g_{x_0}]))
	= h_{*y}([\ol{\alpha_0}\cdot g_{x_0}\cdot \alpha_0])
	= [(h\circ\ol{\alpha_0})\cdot(h\circ g_{x_0})\cdot(h\circ\alpha_0)]
	= \widehat{h\circ\alpha_0}([h\circ g_{x_0}]).
    \end{equation*}
    Here $h\circ\alpha_0$ is a path from $h(x_0) = x_1$ to $h(y_0) = y_1$,
    as is $\alpha_1$. By Munkres's Exercise 52.3 (from HW10),
    $\widehat{\alpha_1} = \widehat{h\circ\alpha_0}$, so
    \begin{equation*}
	h_{*y}(\widehat{\alpha_0}([g_{x_0}]))
	= \widehat{h\circ\alpha_0}([h\circ g_{x_0}])
	= \widehat{\alpha_1}([h\circ g_{x_0}])
	= \widehat{\alpha_1}(h_{*x}(g_{x_0})),
    \end{equation*}
    proving the claim (since $[g_{x_0}]$ generates $\pi_1(S^1,x_0)$).
    If $h_{*x}(\gamma_{x_0}) = d\gamma_{x_1}$, then
    $\widehat{\alpha_1}(h_{*x}(\gamma_{x_0})) = d\gamma_{y_1}$ by definition
    of $\gamma_\bullet$, while
    if $h_{*y}(\gamma_{y_0}) = d'\gamma_{y_1}$, then
    \begin{equation*}
	h_{*y}(\widehat{\alpha_0}(\gamma_{x_0}))
	= h_{*y}(\gamma_{y_0})
	= d'\gamma_{y_1},
    \end{equation*}
    but we proved above that $d\gamma_{y_1} = d'\gamma_{y_1}$,
    so $d = d'$, hence $d$ is independent of the choice of basepoint.
    \\

    (b) Fix $x_0\in S^1$ and let $A\colon S^1\times I\to S^1$ be a homotopy
    such that $A(x,0) = h(x)$ and $A(x,1) = k(x)$.
    If $g_{x_0}\in\gamma_{x_0}$ is a path in the path homotopy class $\gamma_{x_0}$,
    then $B(s,t)\coloneq A(g_{x_0}(s),t)$ is a homotopy (not path homotopy)
    between the paths $h(g_{x_0})$ and $k(g_{x_0})$.

    If we consider $S^1$ as a subset of $\CC$, then
    multiplying $h$ and $k$ by some fixed $z\in S^1$
    does not change their degree, since all it does is change the basepoint
    of the paths when considering the induced homomorphisms.
    Thus, we can assume $h$ and $k$ both send $x_0$ to $1\in\CC$.
    In this case we can modify $B$ to be a \emph{path} homotopy
    between $h(g_{x_0})$ and $k(g_{x_0})$, as
    \begin{equation*}
	C(s,t)\coloneq\frac{B(s,t)}{B(0,t)}
    \end{equation*}
    is still continuous (because the denominator is continuous in $t$
    and is never zero) and is now always equal to $1$ when $s = 0,1$.
    In sum, $h(g_{x_0})$ and $k(g_{x_0})$ are path homotopic,
    so they lie in the same homotopy class
    $\deg(h)\cdot\gamma_1 = \deg(k)\cdot\gamma_1$
    in $\pi_1(S^1,1)$, proving $\deg h = \deg k$.
    \\

    (c) The degree of $h\circ k$
    is the coefficient of $\gamma_{(h\circ k)(x_0)} = \gamma_{h(k(x_0))}$ in
    \begin{equation*}
	(h\circ k)_*(\gamma_{x_0}) = h_*(k_*(\gamma_{x_0}))
	= h_*(\deg k\cdot \gamma_{k(x_0)})
	= \boxed{\deg h\cdot\deg k}\cdot \gamma_{h(k(x_0))}.
    \end{equation*}

    (d) The identity map induces the identity isomorphism
    $\pi_1(S^1,x_0)\to\pi_1(S^1,x_0)$, so its degree is $\boxed 1$
    ($\gamma_{x_0}$ gets sent to itself, $1\cdot \gamma_{x_0}$).

    The constant map sending all points in $S^1$ to some $c$
    induces the trivial homomorphism $\pi_1(S^1,x_0)\to\pi_1(S^1,c)$
    since the image of each path under the constant map is constant.
    Thus, the degree is $\boxed 0$
    ($\gamma_{x_0}$ gets sent to $0\cdot\gamma_{c}$).

    The reflection map sends a path $\alpha$ based at $b_0$ to $\ol\alpha$,
    so the induced homomorphism maps $\gamma_{b_0}$ to $-\gamma_{b_0}$,
    hence its degree is $\boxed{-1}$.

    The map $z\mapsto z^n$ sends loop $\alpha$ based at $b_0 = 1\in\CC$ to
    (a loop homotopic to) $\alpha$ concatenated with itself $n$ times, namely
    \begin{equation*}
	\beta(t) = \alpha(nt - \floor{nt})
    \end{equation*}
    where $\floor x$ is the greatest integer that is at most $x$.
    Thus $h_*([\alpha]) = n[\alpha]$ implies that
    $h_*([\gamma_{b_0}]) = n\gamma_{b_0}$, so the degree of $h$ is $\boxed n$.
\end{soln}

\pagebreak
\begin{problem}
    [Munkres 59.3b]
    Show that $\RR^2$ and $\RR^n$ are not homeomorphic if $n > 2$.
\end{problem}

\begin{soln}
    If $\RR^2$ and $\RR^n$ are homeomorphic
    then the spaces obtained by removing one point from each space
    ought to be homeomorphic as well.
    However, by Munkres's Theorem 58.2,
    $\RR^2$ minus a point has fundamental group isomorphic to that of $S^1$,
    namely $\ZZ$, while for $n>2$, $\RR^2$ minus a point has fundamental group
    isomorphic to that of $S^{n-1}$, which is trivial by Munkres's Theorem 59.3.
    Homeomorphic spaces have the same fundamental group,
    so the inequality of the fundamental groups of $S^1$ and $S^{n-1}$ for $n>2$
    implies that $\RR^2$ and $\RR^n$ cannot be homeomorphic.
\end{soln}

\begin{problem}
    [Munkres 60.1]
    Compute the fundamental groups of the ``solid torus'' $S^1\times B^2$
    and the product space $S^1\times S^2$.
\end{problem}

\begin{soln}
    By Munkres's Theorem 60.1,
    \begin{align*}
        \pi_1(S^1\times B^2, (x_0,y_0)) &\cong
	\pi_1(S^1,x_0)\times\pi_1(B^2,y_0)\cong \ZZ\\
	\pi_1(S^1\times S^2, (x_0,y_0)) &\cong
	\pi_1(S^1,x_0)\times\pi_1(S^2,y_0)\cong \ZZ
    \end{align*}
    since $B^2$ and $S^2$ both have trivial fundamental groups.
\end{soln}

\begin{problem}
    [Munkres 70.1]
    Suppose that the homomorphism $i_*$ induced by the inclusion
    $i\colon U\cap V\to X$ is trivial.
    \begin{enumerate}[(a)]
        \ii Show that $j_1$ and $j_2$ induce an epimorphism
	\begin{equation*}
	    h\colon(\pi_1(U,x_0)/N_1)*(\pi_2(V,x_0)/N_2)\to \pi_1(X,x_0),
	\end{equation*}
	where $N_1$ is the least normal subgroup of $\pi_1(U,x_0)$
	containing the image of $i_1$ and
	$N_2$ is the least normal subgroup of $\pi_1(V,x_0)$
	containing the image of $i_2$.
	\ii Show that $h$ is an isomorphism.
    \end{enumerate}
\end{problem}

\begin{soln}
    Write $U_1$ and $U_2$ instead of $U$ and $V$.

    (a) Let $k\in\{1,2\}$. Since $i_* = j_k\circ i_k$ is trivial,
    $\operatorname{im}i_k\subseteq\ker j_k$, so $N_k\subseteq\ker j_k$.
    Letting $p_k\colon\pi_1(U_k,x_0)\to\pi_1(U_k,x_0)/N_k$
    be the projection homomorphism, the fact that $N_k\subseteq\ker j_k$ means
    there exists a unique homomorphism
    $\ol{j_k}\colon\pi_1(U_k,x_0)/N_k\to\pi_1(X,x_0)$
    such that $\ol{j_k}\circ p_k = j_k$.
    The map
    \begin{equation*}
	h\colon(\pi_1(U_1,x_0)/N_1)*(\pi_2(U_2,x_0)/N_2)\to \pi_1(X,x_0)
    \end{equation*}
    induced by $\ol{j_1}$ and $\ol{j_2}$ has the same image
    as the map $j$ of Munkres's Theorem 70.2
    because $\ol{j_k}$ and $j_k$ have the same image,
    and the aforementioned theorem tells us
    that $j$ is an epimorphism, proving the desired result.
    \\

    (b) Since $h$ is an epimorphism it suffices to show that
    it has a left inverse. Let
    \begin{equation*}
        \Phi\colon\pi_1(X,x_0)\to(\pi_1(U_1,x_0)/N_1)*(\pi_2(U_2,x_0)/N_2)
    \end{equation*}
    be the homomorphism given by taking
    $H = (\pi_1(U_1,x_0)/N_1)*(\pi_2(U_2,x_0)/N_2)$ in van Kampen's theorem
    (the maps $\phi_k\colon\pi_1(U_k,x_0)\to H$ are
    given by projecting $\pi_1(U_k,x_0)$ via $p_k$ and then
    embedding the image $\pi_1(U_k,x_0)/N_k$ into $H$ in the natural way).

    It suffices to check that $\Phi(h(\alpha)) = \alpha$
    for all $\alpha N_k$ in $\pi_1(U_k,x_0)/N_k$
    (which is the coset of $N_k$ by $\alpha\in\pi_1(U_k,x_0)$),
    since the free product is generated by the elements in the factors.
    For such an $\alpha N_k\in\pi_1(U_k,x_0)/N_k$,
    $h(\alpha) = \ol{j_k}(\alpha N_k) = j_k(\alpha)$ by construction.
    The commutative diagram shown in the theorem statement
    says that $\Phi\circ j_k = \phi_k$, so
    \begin{equation*}
        \Phi(h(\alpha N_k)) = \Phi(j_k(\alpha)) = \phi_k(\alpha) = \alpha N_k
    \end{equation*}
    since, as described above, $\phi_k$ simply embeds the coset
    $p_k(\alpha) = \alpha N_k$ into $H$ in the natural way.
    This shows that $\Phi$ is a left inverse of $h$ as desired.
\end{soln}










\end{document}

